% !TeX root = ../main.tex
%%% ---------------------------------------------------------------------------
%%% 选题背景与意义
%%%
%%% 这一节回答委员会的第一个问题：**这件事值得做一个学位吗**。
%%% 目标页数：2~3 页（20 分钟答辩）。讲太久是最常见的时间分配错误——
%%% 背景讲了 8 分钟，自己的工作只剩 6 分钟，这是答辩失分的头号原因。
%%%
%%% 开题答辩：这一节要扩到 4~5 页，因为开题的重点就是论证选题。
%%% 中期/毕业答辩：压到 2 页，委员会已经看过开题了。
%%%
%%% 帧标题写成断言句（见 references/presenting.md 的 Assertion-Evidence）。
%%% ---------------------------------------------------------------------------
\section{选题背景与意义}

\begin{frame}{遥感目标检测的瓶颈在小目标，而非整体精度}
  \begin{columns}[T, onlytextwidth]
    \begin{column}{0.5\textwidth}
      \begin{itemize}
        \setlength{\itemsep}{0.8em}
        \item 对地观测、灾害评估、城市管理的共性技术环节
        \item 现有方法总体 mAP 已达 \qty{75}{\percent} 以上
        \item 但小目标（$<32^2$ 像素）子集只有 \qty{48}{\percent}，
              且贡献了 \qty{78}{\percent} 的漏检
      \end{itemize}
    \end{column}
    \begin{column}{0.46\textwidth}
      \centering
      \placeholder{\linewidth}{4.2cm}{figures/scale-gap.pdf}
    \end{column}
  \end{columns}
  \takeaway{总体指标已经饱和，剩下的空间几乎全在小目标上}

  \note{开场第一页。语速放慢，让委员会进入状态。
        「总体饱和、空间在小目标」这句话要说清楚，它是全篇的立足点。}
\end{frame}

\begin{frame}{应用侧的需求是实时，这限制了可用的解法}
  \begin{evidence}
    \placeholder{0.72\linewidth}{4.4cm}{figures/application-constraints.pdf}
  \end{evidence}
  \takeaway{离线的高精度方案解决不了问题，必须在实时预算内提升小目标}
\end{frame}
